% Evaluación por días: contracción, corrección por viaje y comparación con otros métodos (memoria, capítulo 8).

\subsection{Day-by-day evaluation of the components}
\label{sec:day-by-day}

Rules enter the method only after a day-by-day evaluation that reproduces how they will be used: for each of seven
test days, everything is learned from the previous days only, and the prediction is evaluated end to end on chained
runs of segments (the sum is what the user sees, and a change can improve each segment without improving the sum).
A rule is adopted if it improves at every horizon and on almost every day.

\paragraph{Shrinkage} Table~\ref{tab:shrinkage} compares the band median alone ($n_0=0$) with the shrunk value of
Eq.~\eqref{eq:shrinkage} on 195,786 segment samples over 26 days (test days 25--31 October 2026). With $n_0=5$ every
measure improves at every horizon. With $n_0=10$ the gain holds at short range, but at 12--20 minutes it is worse than
not shrinking: too much shrinkage erases real differences between bands, which accumulate over many segments, as
Proposition~\ref{prop:shrinkage} predicts.

\begin{table}[t]
\centering
\caption{Day-by-day evaluation of the shrinkage towards the segment median: mean absolute error of the sum of
segments and share within $\pm120$~s.}
\label{tab:shrinkage}
\small
\begin{tabular}{@{}lrrrrrrr@{}}
\toprule
& & \multicolumn{3}{c}{Mean error (s)} & \multicolumn{3}{c}{Within $\pm120$~s (\%)} \\
\cmidrule(lr){3-5}\cmidrule(l){6-8}
Horizon & Runs & $n_0=0$ & $n_0=5$ & $n_0=10$ & $n_0=0$ & $n_0=5$ & $n_0=10$ \\
\midrule
4--8 min   & 107,052 & 42.8 & \textbf{42.0} & 42.3 & 96.0 & \textbf{96.2} & 96.0 \\
8--12 min  & 83,336  & 62.7 & \textbf{61.7} & 62.4 & 87.4 & \textbf{87.9} & 87.5 \\
12--20 min & 113,581 & 89.3 & \textbf{88.2} & 89.4 & 73.6 & \textbf{74.1} & 73.4 \\
\bottomrule
\end{tabular}
\end{table}

\paragraph{Per-trip correction} On 87,333 test passages (26 October--1 November 2026), the corrected timetable without
any live data had a mean absolute error of 210~s, against 438~s for the published timetable, and arrived on average
131~s before the bus. With the correction by trip and third of the route, the mean error fell to 158~s and the share
within $\pm2$~minutes rose from 47\% to 54\%, while a single correction per band lowered the mean error but worsened
the tight hit rates, because the bias is not the same for every line and trip. Splitting the route into fifths or stops
did not help. Replacing the mean of the deviations by their 40th percentile, evaluated by reconstructing what the app
would have shown on two days not seen during learning, raised the share within $\pm2$~minutes by 4.1 and 2.5 points
and cut by a third the cases in which the bus passed more than two minutes before the time shown.

\paragraph{Comparison with other methods} Table~\ref{tab:methods} compares the method, under the same protocol and on
the same runs, with alternatives that, unlike it, also use what is happening that day: a single scale factor; the
recent ratio of the last buses on each segment; a Kalman filter per segment in the spirit of \textcite{shalaby2004};
a Kalman filter on the vehicle's own pace; and a neural network with two hidden layers of 32 and 16 units and 28
inputs, in the spirit of \textcite{chien2002} and \textcite{jeong2005}, trained for each test day on the previous days.
Hyper-parameters were tuned on the seven days before the test period. The method keeps almost all the gain over the
timetable: at 4--8 minutes the mean error falls from 154~s to 42~s. The alternatives improve on it a little and on
almost every day: buses that have just run the same segments, alone or through a Kalman filter, by 1--2\%; the
vehicle's own pace by about 4\%; and the neural network, combining everything, by 6--7\%, the only one better on all
seven days at all three horizons (95\% bootstrap intervals of the difference exclude zero). These comparisons are on
sums of segments, not on the full arrival shown to the user.

\begin{table}[t]
\centering
\caption{Comparison with other methods under the day-by-day protocol (test days 25--31 October 2026): mean absolute
error of the sum of segments, share within $\pm120$~s at 12--20 minutes, and number of test days (of seven) on which
the alternative beats the method, at each horizon.}
\label{tab:methods}
\footnotesize
\setlength{\tabcolsep}{4pt}
\begin{tabular}{@{}>{\raggedright\arraybackslash}p{3.2cm}rrrrc@{}}
\toprule
& \multicolumn{3}{c}{Mean error (s)} & Within & Days \\
\cmidrule(lr){2-4}
Method & 4--8 min & 8--12 min & 12--20 min & $\pm120$~s (\%) & better \\
\midrule
Published timetable & 153.9 & 233.5 & 329.7 & 24.0 & 0 \\
\textbf{This method} & \textbf{42.0} & \textbf{61.7} & \textbf{88.2} & \textbf{74.1} & --- \\
Scaled by one factor & 42.6 & 61.9 & 86.9 & 74.9 & 2 · 4 · 5 \\
Previous buses & 41.6 & 60.7 & 86.5 & 75.0 & 5 · 7 · 7 \\
Kalman filter per segment & 41.7 & 60.7 & 86.4 & 75.0 & 5 · 7 · 7 \\
Kalman filter on own pace & 40.4 & 59.3 & 84.6 & 76.0 & 6 · 7 · 7 \\
Neural network, no same-day data & 40.7 & 58.8 & 82.5 & 77.2 & 6 · 6 · 6 \\
Neural network & 39.5 & 57.5 & 81.7 & 77.4 & 7 · 7 · 7 \\
\bottomrule
\end{tabular}
\end{table}
